\documentclass[10pt]{article}
\usepackage[T1]{fontenc}
\usepackage[utf8]{inputenc}
\usepackage[french]{babel}
\usepackage{amsmath,amssymb,mathrsfs,mathpazo}
\usepackage{geometry}
\geometry{a4paper,margin=2cm}
\usepackage{graphicx,array,tabularx,tabularray,colortbl,multirow,multicol,booktabs,enumitem,pifont,ulem,fancybox,diagbox,cancel}
\usepackage{pstricks,pstricks-add,pst-plot,pst-tree,pst-node,pst-text,pst-func,pst-math,pst-bspline,pst-3dplot,pst-eucl,pst-solides3d}
\let\pscalculate\undefined
\usepackage{pst-bezier}
\usepackage{tikz,pgfplots,tkz-tab}
\usetikzlibrary{arrows,arrows.meta,patterns,calc,positioning,trees,decorations.markings,shapes.geometric,shapes.misc,shapes.arrows,intersections,decorations.pathmorphing,decorations.pathreplacing,backgrounds,shadows,babel}
\pgfplotsset{compat=1.18}
\usepgfplotslibrary{fillbetween}
\usepackage[np]{numprint}
\usepackage{listings,makecell,pas-tableur,twoopt}
\let\barre\undefined
\usepackage[pstricks]{tablvar}
\definecolor{darkgreen}{RGB}{0,100,0}
\providecommand{\red}{\color{red}}
\providecommand{\blue}{\color{blue}}
\providecommand{\green}{\color{green}}
\providecommand{\black}{\color{black}}
\providecommand{\white}{\color{white}}
\providecommand{\gray}{\color{gray}}
\providecommand{\cyan}{\color{cyan}}
\providecommand{\magenta}{\color{magenta}}
\providecommand{\yellow}{\color{yellow}}
\providecommand{\violet}{\color{violet}}
\providecommand{\orange}{\color{orange}}
\providecommand{\euro}{\texteuro}
\providecommand{\EUR}{\texteuro}
\providecommand{\ieme}{\textsuperscript{e}}
\providecommand{\degres}{\ensuremath{^\circ}}
\providecommand{\degree}{\ensuremath{^\circ}}
\providecommand{\vv}[1]{\overrightarrow{#1}}
\providecommand{\Coord}[1]{#1}
\providecommand{\decosix}{}
\providecommand{\starredbullet}{\ensuremath{\bullet}}
\pagestyle{empty}
\setlength{\parindent}{0pt}
\usepackage{adjustbox}\DeclareUnicodeCharacter{2236}{:}\DeclareUnicodeCharacter{2190}{\ensuremath{\leftarrow}}
\providecommand{\R}{\mathbb{R}}
\providecommand{\N}{\mathbb{N}}
\providecommand{\Z}{\mathbb{Z}}
\providecommand{\Q}{\mathbb{Q}}
\providecommand{\D}{\mathbb{D}}
\providecommand{\C}{\mathbb{C}}
\providecommand{\e}{{\mathrm{e}}}
\providecommand{\vect}[1]{\overrightarrow{#1}}
\providecommand{\vectt}[1]{\overrightarrow{\mathrm{#1}}}
\providecommand{\vvt}[1]{\overrightarrow{\mathrm{#1}}}
\providecommand{\np}[2][]{{#2~#1}}
\providecommand{\vv}[1]{{\overrightarrow{#1}}}
\renewcommand{\d}{{\mathrm{d}}}
\providecommand{\degres}{{^\circ}}
\providecommand{\degre}{{^\circ}}
\providecommand{\blue}{\color{blue}}
\providecommand{\red}{\color{red}}
\providecommand{\enskip}{\quad}
\providecommand{\llbracket}{\lbrack\!\lbrack}
\providecommand{\rrbracket}{\rbrack\!\rbrack}
\providecommand{\euro}{{\text{€}}}
\providecommand{\no}{{\text{n°}}}
\providecommand{\interval}[2]{{[#1;#2]}}
\providecommand{\intervFO}[2]{{[#1;#2[}}
\providecommand{\dotfill}{\ldots}

\providecommand{\ds}{\displaystyle}
\providecommand{\textding}[1]{#1}
\providecommand{\Oij}{\(\mathrm{O};\vec{\imath},\vec{\jmath}\)}
\providecommand{\Oijk}{\(\mathrm{O};\vec{\imath},\vec{\jmath},\vec{k}\)}
\providecommand{\Ouv}{\(\mathrm{O};\vec{u},\vec{v}\)}
\providecommand{\intervFF}[2]{\ensuremath{[#1;#2]}}
\renewcommand{\euro}{\ifmmode\text{\texteuro}\else\texteuro\fi}
\DeclareUnicodeCharacter{00B0}{\ifmmode^{\circ}\else\textdegree\fi}
\DeclareUnicodeCharacter{20AC}{\ifmmode\text{\texteuro}\else\texteuro\fi}
% Énoncés : Éducation nationale. Transcription et corrigés : APMEP, Denis Vergès et contributeurs. Découpage et classement : CoopMaths, Sébastien Lozano, Éric Elter et contributeurs. Adaptation de la présentation : Maths Explorer.
% CC BY-SA : https://coopmaths.fr/www/about/
% https://coopmaths.fr/alea/static/e3c/2020/tex/e3c_2020_00_sujet49_1.tex
\begin{document}
\noindent\textit{APMEP / CoopMaths. CC BY-SA. Présentation adaptée par Maths Explorer.}\par\medskip
\medskip

Ce QCM comprend 5 questions indépendantes.

Pour chacune d'elles, une seule des affirmations proposées est exacte.
 
Indiquer pour chaque question sur la copie la lettre correspondant à la réponse choisie.

 Aucune justification n'est demandée.
 
Chaque réponse correcte rapporte 1 point. Une réponse incorrecte ou une absence de réponse n'apporte ni ne retire de point.

\begin{enumerate}
\item On considère la droite $d$ dont une équation cartésienne dans un repère orthonormé est $2x-3y+4=0$.

\begin{tabularx}{\linewidth}{*{2}{X}}
\textbf{a.~~} Un vecteur directeur de $d$ est $\vv{u} \binom{-6}{4}$. &\textbf{b.~~} Un vecteur normal de $d$ est $\vv{n}\binom{-12}{18}$. \\
\textbf{c.~~} Le point $ C(-5~;~2)$ appartient à la droite $d$.& \textbf{d.~~} La droite $d$ coupe la droite d'équation 

$-x+3y-2=0$ au point F (1~;~2).
\end{tabularx}

\item Dans un repère orthonormé le cercle $\mathcal{C}$ a pour équation $x^2-2x+y^2+y=3$ 

et la droite $\mathcal{D}$ pour équation $y=1$.


\begin{tabularx}{\linewidth}{*{1}{X}}
\textbf{a.~~} $\mathcal{C}$ et $\mathcal{D}$ n'ont aucun point d'intersection.\\
\textbf{b.~~} $\mathcal{C}$ et $\mathcal{D}$ ont un seul point d'intersection.\\
\textbf{c.~~} $\mathcal{C}$ et $\mathcal{D}$ ont deux points d'intersection.\\
\textbf{d.~~}On ne peut pas savoir combien $\mathcal{C}$ et $\mathcal{D}$ ont de points d'intersection.\\
\end{tabularx}

\item La fonction $f$ est la fonction définie sur l'ensemble des réels par $f(x)= \cos(2x)$.

\begin{tabularx}{\linewidth}{*{2}{X}}
\textbf{a.~~}$f $est paire.  &\textbf{b.~~}$f$ est impaire. \\\textbf{c.~~}$f$ n'est ni paire ni impaire.& \textbf{d.~~}$f$ a pour période $\dfrac{\pi}{2}$ .
\end{tabularx}

\item Soit la suite $\left(u_n\right)$ définie par $u_0=1$ et pour tout entier naturel $n$ par $u_{n+1}=\dfrac{1}{2}\left(u_n + \dfrac{2}{u_n}\right)$.

On définit en langage Python une fonction \og suite \fg{} pour calculer $u_n$ connaissant $n$.

\begin{tabular}{*{2}{lr}}
\textbf{a.~~} 
\begin{minipage}[]{0.4\linewidth}
\begin{lstlisting}[style=python]
 def suite(n):
	u=0 
	for i in range (1,n+1):
		u=1/2*(u+2/u) 
	return u
 \end{lstlisting}
\end{minipage} 
 &
\textbf{b.~~}
\begin{minipage}[]{0.4\linewidth}
\begin{lstlisting}[style=python]
 def suite(n):
	u=1 
	for i in range (1,n+1):
		u=1/2*(u+2/u)
	return n
\end{lstlisting}
\end{minipage}\\

\textbf{c.~~}
\begin{minipage}[]{0.4\linewidth}
\begin{lstlisting}[style=python]
 def suite(n):
	u=1
	for i in range (1,n+1):
		u=1/2*u+2/u 
	return u
\end{lstlisting}
\end{minipage}&
\textbf{d.~~}
\begin{minipage}[]{0.4\linewidth}
\begin{lstlisting}[style=python]
 def suite(n):
	u=1 
	for i in range (1,n+1):
		u=1/2*(u+2/u)
	return u 
\end{lstlisting}
\end{minipage}
\end{tabular}

\item L'équation $\e^x=1$ :

\begin{tabularx}{\linewidth}{*{2}{X}}
\textbf{a.~~}n'a pas de solution.  &\textbf{b.~~}a pour solution le nombre 1. \\\textbf{c.~~}a pour solution le nombre 0.& \textbf{d.~~} a pour solution le nombre $\e$.
\end{tabularx}
\end{enumerate}

\vspace{0,5cm}

\end{document}
